\section{The support functional and its maximum}\label{sec:calibration}

Put $L=\pi/2$. This section concerns pairs of $C^1$ functions on
$[0,L]$, without a convexity assumption. For such a pair $X=(F,G)$ let
\[
 p_X=F'-G+1,\qquad q_X=G'+F-1,
 \qquad p_-:=\min(p,0),\quad q_+:=\max(q,0),
\]
and define
\begin{equation}\label{eq:support-functional}
 \begin{split}
 \mathcal A(X)=\frac12\int_0^L\bigl[&F^2+G^2-F'^2-G'^2
                  +(F-1)^2+(G-1)^2\\
       &+(F-1)G'-(G-1)F'-p_{X,-}^2-q_{X,+}^2\bigr]\,\dd t.
 \end{split}
\end{equation}
The admissible traces at half-width $a$ are
\begin{equation}\label{eq:calibration-traces}
 F(0)=G(L)=a,\qquad F(L)=G(0)=1.
\end{equation}
The support-area identity and the three niche contact arcs give this
functional: their common corner terms cancel, leaving the two squared
contact momenta as penalties. We can maximize it on a larger class of
profiles than the convex supports. At fixed traces a stationary pair
will be the unique maximum; varying the common side trace will then
select the reference width.

\subsection{The fixed-trace remainder}

The positive-part squares already have nonnegative convex remainders.
The remaining quadratic form has a mixed derivative term. Rotating the
pair $(\phi,\psi)$ through angle $-t/2$ absorbs that term into a
derivative square. A Dirichlet estimate then makes the remainder
strictly positive for every nonzero variation with fixed traces.

For a variation $V=(\phi,\psi)$ write
\[
 B_0(V)=\frac12\int_0^L
 \bigl[\phi'^2+\psi'^2-2\phi^2-2\psi^2
                      -\phi\psi'+\psi\phi'\bigr]\,\dd t.
\]
The first variation of \cref{eq:support-functional} is
\begin{align}\label{eq:calibration-first-variation}
 D\mathcal A(X)[V]=\int_0^L\bigl[&
 (2F-1)\phi+(2G-1)\psi-F'\phi'-G'\psi'\notag\\
 &+\tfrac12\bigl(\phi G'+(F-1)\psi'
                       -\psi F'-(G-1)\phi'\bigr)\notag\\
 &-p_{X,-}(\phi'-\psi)-q_{X,+}(\psi'+\phi)\bigr]\,\dd t.
\end{align}
These formulas make sense when the first derivatives are continuous;
the stationary pairs constructed below have piecewise continuous
second derivatives.

\begin{lemma}[Dirichlet estimate]\label{lem:calibration-dirichlet}
If $f\in C^1[0,L]$ and $f(0)=f(L)=0$, then
\[
 \int_0^L f^2\le\frac{L^2}{6}\int_0^L f'^2
\]
\end{lemma}
\begin{proof}
Cauchy--Schwarz applied on $[0,t]$ and $[t,L]$ gives
\[
 f(t)^2\le t\int_0^t f'^2,\qquad
 f(t)^2\le(L-t)\int_t^L f'^2.
\]
Multiply the first inequality by $L-t$, the second by $t$, and add.
The result is
$Lf(t)^2\le t(L-t)\int_0^L f'^2$. Integration gives the
inequality, since $\int_0^L t(L-t)\,\dd t=L^3/6$.
\end{proof}

\begin{lemma}[Coercive remainder]\label{lem:calibration-coercivity}
For a variation $V$ vanishing at all four traces, put
$y(t)=e^{-it/2}(\phi(t)+i\psi(t))$. Then
\begin{equation}\label{eq:calibration-coercivity}
 B_0(V)=\frac12\int_0^L
           (|y'|^2-\tfrac94|y|^2)\,\dd t
       \ge\frac12\left(1-\frac{3L^2}{8}\right)
             \int_0^L|y'|^2\,\dd t.
\end{equation}
Equality $B_0(V)=0$ implies $V=0$ on $[0,L]$.
\end{lemma}
\begin{proof}
For $z=\phi+i\psi$,
$|y'|^2=|z'|^2+|z|^2/4-\operatorname{Im}(\overline z z')$.
This gives the identity in \cref{eq:calibration-coercivity}.
Apply \cref{lem:calibration-dirichlet} to the two real components of
$y$. The coefficient is positive because
$3L^2/8=3\pi^2/32<1$.
If $B_0=0$, the derivative integral is zero, so $y$ is constant.
Its zero traces imply $y=0$ and hence $V=0$.
\end{proof}

\begin{theorem}[Fixed-trace maximum]\label{thm:fixed-trace-max}
Suppose $X$ satisfies \cref{eq:calibration-traces} and
$D\mathcal A(X)[V]=0$ for every zero-trace $C^1$ variation $V$.
For every pair $\widetilde X$ with the same traces,
\begin{equation}\label{eq:fixed-trace-max}
 \mathcal A(\widetilde X)\le\mathcal A(X)-B_0(\widetilde X-X).
\end{equation}
Equality of the two functional values implies $\widetilde X=X$.
\end{theorem}
\begin{proof}
The function $s\mapsto s_-^2$ is convex and differentiable, with
derivative $2s_-$. Its remainder
\[
 R_-(s,t)=t_-^2-s_-^2-2s_-(t-s)
\]
is nonnegative. Define $R_+(s,t)$ similarly with positive parts.
Expanding the polynomial terms in
\cref{eq:support-functional} gives the exact identity
\begin{align*}
 \mathcal A(\widetilde X)=\mathcal A(X)
 &+D\mathcal A(X)[\widetilde X-X]-B_0(\widetilde X-X)\\
 &-\tfrac12\int_0^L
  \bigl(R_-(p_X,p_{\widetilde X})+R_+(q_X,q_{\widetilde X})\bigr)\,\dd t.
\end{align*}
Stationarity and nonnegativity give \cref{eq:fixed-trace-max}.
If the values agree, \cref{lem:calibration-coercivity} makes the
variation zero pointwise.
\end{proof}

The fluxes and forces associated with the first variation are
\begin{align*}
 J_F&=-F'-(G-1)/2-p_-,& K_F&=2F-1+G'/2-q_+,\\
 J_G&=-G'+(F-1)/2-q_+,& K_G&=2G-1-F'/2+p_-.
\end{align*}
Thus $D\mathcal A(X)[V]=\int(K_F\phi+K_G\psi+J_F\phi'+J_G\psi')$.
If $J_F'=K_F$ and $J_G'=K_G$, integration by parts gives
\begin{equation}\label{eq:calibration-boundary-flux}
 D\mathcal A(X)[V]=[J_F\phi+J_G\psi]_0^L.
\end{equation}
The equations can hold separately on finitely many subintervals,
provided the fluxes are continuous at their endpoints.
For a symmetric pair $G(t)=F(L-t)$ one has
$J_G(t)=-J_F(L-t)$ and $K_G(t)=K_F(L-t)$.
It therefore suffices to verify the first flux equation.

\subsection{The stationary family}

The stationary profiles come from the Euler equations on the three
contact-sign regions. Substitution into $J_F'=K_F$, $J_G'=K_G$ gives
\[
 \begin{array}{c|cc}
 \text{sign region}&\text{first equation}&\text{second equation}\\ \hline
 p,q\ge0&F''+F=0&G''+G=1/2\\
 p\le0\le q&2F''-G'+F=0&2G''+F'+G=0\\
 p,q\le0&F''+F=1/2&G''+G=0.
 \end{array}
\]
Choose the first switch at $b$ and impose reflection symmetry about
$L/2$. The last switch is then $L-b$. Solving these linear equations
and matching values and derivatives gives the coefficients below; the
half-angle term belongs to the coupled middle equations.

For $b\in[0,\pi/4]$ put
\begin{equation}\label{eq:stationary-family-parameters}
 \begin{gathered}
 \gamma=b/2+\pi/8,\qquad k=\cot\gamma,\qquad
 a=\frac{\sin b+2k\cos b}{3},\\
 E=\frac12-\frac{\cos b}{2}+k\sin b,\qquad
 R=\frac{2}{3\sin\gamma},\qquad
 S=\frac12-\frac{\cos b}{6}+\frac{k\sin b}{3}.
 \end{gathered}
\end{equation}
These letters denote functions of $b$ in this subsection.
Define
\begin{equation}\label{eq:stationary-family}
 F_b(t)=
 \begin{cases}
 a\cos t+E\sin t,&0\le t\le b,\\
 R\cos(t/2+\pi/8)+S\sin t,&b\le t\le L-b,\\
 \tfrac12+\tfrac12\sin t+\tfrac a2\cos t,&L-b\le t\le L,
 \end{cases}
 \qquad G_b(t)=F_b(L-t).
\end{equation}

\begin{lemma}[Matching and stationarity]\label{lem:stationary-family}
The pair $X_b=(F_b,G_b)$ is $C^1$, satisfies
\cref{eq:calibration-traces} at width $a(b)$, and is stationary.
Its contact signs are $p,q\ge0$ on $[0,b]$, $p\le0\le q$ on
$[b,L-b]$, and $p,q\le0$ on $[L-b,L]$.
\end{lemma}
\begin{proof}
The identities $R\sin\gamma=2/3$ and $R\cos\gamma=2k/3$
give value and derivative matching at $b$ by substitution in
\cref{eq:stationary-family}. Reflection gives matching at $L-b$.
The axis values are $F_b(0)=a$ and $F_b(L)=1$.

On the first phase, direct differentiation gives
\[
 p=\tfrac12+(E-\tfrac12)\cos t-\tfrac32a\sin t,
 \qquad q=-1+\tfrac32a\cos t+(E-\tfrac12)\sin t.
\]
Substituting \cref{eq:stationary-family-parameters}, these become
\[
 p=\frac{1-\cos(b-t)}2+k\sin(b-t),\qquad
 q=-1+k\cos(b-t)+\tfrac12\sin(b-t).
\]
They are nonnegative: $k\ge1$, $0\le b-t\le\pi/4$, and
$\cos r+\tfrac12\sin r\ge1$ for $0\le r\le\pi/4$.
For the last inequality, square the nonnegative sides, or use
$\sin r\le\cos r$ and $\sin r\le1$ in
$\cos^2r+\sin^2r=1$.
On the middle phase, put $\varphi=t/2+\pi/8$.
Then $p=1-3R\sin\varphi/2$ and $q=3R\cos\varphi/2-1$.
Since $\gamma\le\varphi\le\pi/2-\gamma$, both sine and cosine
are at least $\sin\gamma$. This proves the middle signs.
Reflection, namely $p(t)=-q(L-t)$, proves the last-phase signs.

The outer phases satisfy the uncoupled Euler equations by their
sine and cosine form. In the middle, $G_b=R\sin\varphi+S\cos t$;
the half-angle terms cancel in $2F_b''-G_b'+F_b$, and the
full-angle terms cancel separately. Reflection gives the other
equation. The proved signs select these Euler equations, and
$C^1$ matching makes the fluxes continuous at the switches.
Thus \cref{eq:calibration-boundary-flux} gives stationarity.
\end{proof}

\begin{lemma}[Value of the family]\label{lem:stationary-family-value}
At the width $a(b)$,
\begin{equation}\label{eq:stationary-family-score}
 \mathcal A(X_b)-a(b)=
 H(b):=\frac{(2E-1)(a+\sin b)+3a\cos b+b-2}{2}.
\end{equation}
\end{lemma}
\begin{proof}
Subtracting the first variation in the direction $X$ from twice
the functional cancels its quadratic terms. Write $A_X(t)$ for
the density in \cref{eq:support-functional} and $D_X(t)$ for
that first-variation density.
By \cref{eq:calibration-boundary-flux},
\[
 \int_0^L D_X=2\bigl(J_F(L)F(L)-J_F(0)F(0)\bigr).
\]
On the first phase $2A_X-D_X$ has primitive
\[
 V_E(t)=t/2+(E/2+a/4-1/2)\sin t+(a+1/4)\cos t.
\]
On the last phase its primitive is
\[
 V_W(t)=t/2-(E/2+a/4-1/2)\cos t-(a+1/4)\sin t.
\]
On the middle phase $2A_X-D_X=(G'-F')/2$; this follows by
substituting $p_-=p$ and $q_+=q$ and cancelling the quadratic terms.
Integrating the three primitives and adding the boundary term,
using $F(0)=a$, $F(L)=1$, gives
\[
 2\mathcal A(X_b)=(2E-1)(a+\sin b)+3a\cos b+b-2+2a.
\]
This is \cref{eq:stationary-family-score}.
\end{proof}

\begin{lemma}[The unique scalar maximum]\label{lem:calibration-scalar}
The width $a(b)$ is strictly decreasing from
$a(0)=\tfrac23\cot(\pi/8)$ to $a(\pi/4)=\sqrt2/2$.
The function $H$ has a unique maximum at
$\beta=\arctan Y$, and $H(\beta)=\Mc/2$.
\end{lemma}
\begin{proof}
Differentiating $k$ gives $k'=-(1+k^2)/2$. Hence
\[
 a'=-\frac{k(k\cos b+2\sin b)}3<0,
 \qquad (2E-1)'=k(2\cos b-k\sin b)>0.
\]
The second sign uses $k\sin b\le\cos b$, which follows from
$\gamma\ge b$ and the sine difference formula.
The derivative of the score follows from the boundary flux.
Differentiate the three phase integrals defining $\mathcal A(X_b)$.
Their moving-endpoint terms cancel because the profiles and their
first derivatives match. The Euler equations remove the interior
terms, leaving \cref{eq:calibration-boundary-flux}.
Only $F_b(0)=G_b(L)=a(b)$ vary, and
$J_F(0)=-E$, $J_G(L)=E$. Consequently
\[
 \frac{\dd}{\dd b}\mathcal A(X_b)=2E(b)a'(b),\qquad
 H'(b)=(2E(b)-1)a'(b).
\]
The optimal member therefore has endpoint height $E=1/2$:
its area derivative balances the width charge. This gives
$2\cot\gamma\sin b=\cos b$.
To identify its root, set $y=\tan b$. It gives
$\tan\gamma=2y$. Since $2\gamma=b+\pi/4$,
the tangent double-angle formula gives
\[
 \frac{4y}{1-4y^2}=\frac{1+y}{1-y},\qquad
 4y^3+3y-1=0.
\]
The zero exists uniquely because $2E-1$ is strictly increasing
and has opposite signs at $0$ and $\pi/4$; at that zero
$\tan\gamma=2y<1$, so the denominators used above are positive.
Thus it is $\beta=\arctan Y$.
The signs of $a'$ and $2E-1$ make $H$ strictly increase before
$\beta$ and strictly decrease after it.
At $\beta$, the equation for $k$ gives
$a=1/(3\sin\beta)$, and therefore
\[
 H(\beta)=\frac{\cot\beta+\beta-2}{2}
       =\frac{1/Y+\beta-2}{2}=\Mc/2.
\]
The last equality uses $1/Y=4Y^2+3$.
\end{proof}

\subsection{The large-width tail}

Put $\gamma_0=\pi/8$. For $r\sin\gamma_0\ge2/3$, define
\[
 F_r(t)=r\cos(t/2+\gamma_0)+(1-r\sin\gamma_0)\sin t,
 \qquad G_r(t)=F_r(L-t).
\]
Its half-width is $a=r\cos\gamma_0$.

\begin{lemma}[Tail comparison]\label{lem:calibration-tail}
The tail pair is stationary at its traces, and
\[
 \mathcal A(F_r,G_r)-a
       =3r\cos\gamma_0-1-3r^2\sin\gamma_0\cos\gamma_0<2/3.
\]
\end{lemma}
\begin{proof}
The momenta are $p=1-3r\sin(t/2+\gamma_0)/2\le0$ and
$q=3r\cos(t/2+\gamma_0)/2-1\ge0$. The middle-phase flux
calculation of \cref{lem:stationary-family} applies on the whole
interval and proves stationarity.
Here $2A_X-D_X=(G'-F')/2$ throughout. Integrating this identity
and the boundary flux gives
$\mathcal A=4r\cos\gamma_0-1-3r^2\sin\gamma_0\cos\gamma_0$.
For $x=r\sin\gamma_0\ge2/3$, one has $x-x^2\le2/9$, so
the score is at most $\tfrac23\cot\gamma_0-1$.
Since $\cot(\pi/8)=1+\sqrt2<5/2$, this is strictly less than $2/3$.
\end{proof}

\begin{theorem}[Global calibration]\label{thm:global-calibration}
If $X=(F,G)\in C^1[0,L]^2$ satisfies
\cref{eq:calibration-traces} with $a\ge4/5$, then
\begin{equation}\label{eq:global-calibration}
 \mathcal A(X)-a\le\Mc/2.
\end{equation}
Equality holds precisely for $a=1/(3\sin\beta)$ and $X=X_\beta$.
\end{theorem}
\begin{proof}
If $a\le a(0)$, \cref{lem:calibration-scalar} and the intermediate
value theorem give a unique $b\in[0,\pi/4]$ with $a(b)=a$,
since $\sqrt2/2<4/5$. Apply \cref{thm:fixed-trace-max} to the
stationary pair $X_b$, followed by \cref{lem:calibration-scalar}.
Equality forces $b=\beta$ and then $X=X_\beta$ pointwise.
If $a\ge a(0)$, choose $r=a/\cos\gamma_0$ and apply the same
fixed-trace theorem to the tail. Its score is less than $2/3<\Mc/2$.
The common boundary $a=a(0)$ is covered by either construction.
\end{proof}

\begin{corollary}[Translation mode]\label{cor:calibration-translation}
Suppose $F(L)=G(0)=1$ and $F(0)+G(L)\ge8/5$. Then
\[
 \mathcal A(F,G)-\frac{F(0)+G(L)}2\le\Mc/2.
\]
Equality determines the reference support pair up to one horizontal
translation.
\end{corollary}
\begin{proof}
The transformation $F_t(s)=F(s)+t\cos s$,
$G_t(s)=G(s)-t\sin s$ preserves $p$ and $q$. Expanding the density
shows that its change is $t$ times the derivative of
\[
 \frac{3(F(s)\sin s+G(s)\cos s)-\cos s-\sin s}{2}.
\]
Hence $\mathcal A(F_t,G_t)=\mathcal A(F,G)$ when $F(L)=G(0)$.
Choose $t=(G(L)-F(0))/2$; the new side traces both equal
$a=(F(0)+G(L))/2$. Apply \cref{thm:global-calibration}.
The centered reference pair is exactly $X_\beta$, as substitution
in its three support phases shows. Undoing the transformation
gives the claimed translation mode.
\end{proof}

\subsection{Recovering every full maximum}

\begin{theorem}[Sharp full-angle cap bound]\label{thm:full-cap-bound}
Every downward cap $U$ satisfies
\begin{equation}\label{eq:full-cap-bound}
 P_L(U)\le\Mc/2.
\end{equation}
Equality holds if and only if $U$ is a horizontal translate of
the reference cap $\KR$.
\end{theorem}
\begin{proof}
By \cref{lem:selection-maximum}, a full maximum exists.
Canonicalization in \cref{lem:selection-canonical} preserves
its maximizing value. Its middle is horizontal by
\cref{lem:full-affine-rigidity}; hence its roof is one on $J_U$.
The horizontal estimates of \cref{lem:full-horizontal-unit}
give, after horizontal translation, support traces
$F(0)=G(L)=a=w_U/2>1$, $F(L)=G(0)=1$, both regular curvature
densities at most one, top face exactly $J_U$, and ordered
contact switches. The support-area and signed-niche identities
of \cref{cor:horizontal-score-identity} therefore give
\[
 P_L(U)=\mathcal A(F,G)-a\le\Mc/2
\]
by \cref{thm:global-calibration}. Since the reference attains
$\Mc/2$, this is the global maximum and proves
\cref{eq:full-cap-bound} for every cap.

The same reasoning applies to any prescribed canonical maximum:
\cref{thm:selection-prescribed} and its limiting laws target that
cap itself. Equality in \cref{thm:global-calibration} forces
$a=1/(3\sin\beta)$ and $(F,G)=X_\beta$. By the reference
support formulas, this is the centered reference cap. A downward
cap is determined by its supports on the upper semicircle:
the cap is the intersection of their supporting halfplanes with
$y\ge0$, and every lower support is already determined by its
baseline. Thus every canonical maximum is a translate of $\KR$.

For an arbitrary equality cap, let $V$ be its middle-chord cut
and let $W=V+[0,\varepsilon]e_y$ be its height-one extrusion,
as in \cref{lem:selection-canonical}. The proved bound makes
all inequalities in
\[
 \Mc/2=P_L(U)\le P_L(V)\le P_L(W)\le\Mc/2
\]
equalities. The canonical classification gives $W=\KR+(b,0)$.
The exact extrusion deficit of \cref{lem:full-extrusion-deficit}
then forces $\varepsilon=0$, since the reference niche is
continuous and zero at both endpoints of its middle window.
Consequently $V=W$ and its roof is one throughout $J_U$.
Inclusion $V\subset U\subset\R\times[0,1]$ makes the original
roof one there too. Outside $J_U$ the chord cut preserved that
roof. Hence $U=V=W$. Translation invariance and the reference
value prove the converse.
\end{proof}

\begin{lemma}[Extrusion deficit]\label{lem:full-extrusion-deficit}
Let $V$ be a positive-width cap and suppose
$W=V+[0,\varepsilon]e_y$ has height at most one. Then
\begin{equation}\label{eq:full-extrusion-deficit}
 N_L(V)=(N_L(W)-\varepsilon)_+,\qquad
 P_L(W)-P_L(V)=\int_{J_V}(\varepsilon-N_L(W)(x))_+\,\dd x.
\end{equation}
If $N_L(W)$ is continuous and vanishes at an endpoint of $J_V$,
the last integral is positive whenever $\varepsilon>0$.
\end{lemma}
\begin{proof}
Vertical Minkowski addition raises the roof by $\varepsilon$
and every upward support by $\varepsilon$ times its vertical
normal coordinate. Both raw inner walls therefore rise by
exactly $\varepsilon$. If $r$ is their signed angular supremum
before extrusion, then $N_L(V)=r_+$ and
$N_L(W)=(r+\varepsilon)_+$. This proves the first identity,
including points where $r$ is negative.
The exterior roof gain is $\varepsilon|I_V\setminus J_V|$;
this equals $\varepsilon|J_V|$. Subtract the niche gain
$N_L(W)-(N_L(W)-\varepsilon)_+$ to obtain the second identity.
Continuity next to a zero endpoint gives an interval of positive
length on which $N_L(W)<\varepsilon/2$, proving strictness.
\end{proof}
